\hypertarget{chapter-12-functional-programming}{%
\section{CHAPTER 12: FUNCTIONAL
PROGRAMMING}\label{chapter-12-functional-programming}}

\hypertarget{functional-programming-in-c11}{%
\section{FUNCTIONAL PROGRAMMING IN
C++11}\label{functional-programming-in-c11}}

\hypertarget{lambda-expressions}{%
\subsection{1. Lambda Expressions}\label{lambda-expressions}}

Anonymous functions defined inline.

\hypertarget{basic-syntax}{%
\subsubsection{1.1 Basic Syntax}\label{basic-syntax}}

\passthrough{\lstinline![captures](params) -> return\_type \{ body \}!}

\begin{lstlisting}[language={C++}]
auto add = [](int a, int b) { return a + b; };
int result = add(2, 3);
\end{lstlisting}

\hypertarget{capture-lists}{%
\subsubsection{1.2 Capture Lists}\label{capture-lists}}

Controls access to outer scope variables.

\begin{itemize}
\tightlist
\item
  \passthrough{\lstinline![]!}: No capture.
\item
  \passthrough{\lstinline![x]!}: Capture \passthrough{\lstinline!x!} by
  value (copy).
\item
  \passthrough{\lstinline![\&x]!}: Capture \passthrough{\lstinline!x!}
  by reference.
\item
  \passthrough{\lstinline![=]!}: Capture all by value.
\item
  \passthrough{\lstinline![\&]!}: Capture all by reference.
\item
  \passthrough{\lstinline![this]!}: Capture class members.
\end{itemize}

\begin{lstlisting}[language={C++}]
int x = 10;
auto addX = [x](int y) { return x + y; }; // x is read-only inside
\end{lstlisting}

\hypertarget{mutable-lambdas}{%
\subsubsection{1.3 Mutable Lambdas}\label{mutable-lambdas}}

By default, value captures are \passthrough{\lstinline!const!}.
\passthrough{\lstinline!mutable!} allows modification.

\begin{lstlisting}[language={C++}]
int x = 0;
auto increment = [x]() mutable { return ++x; }; // Modifies local copy
\end{lstlisting}

\begin{center}\rule{0.5\linewidth}{0.5pt}\end{center}

\hypertarget{stdfunction}{%
\subsection{\texorpdfstring{2.
\texttt{std::function}}{2. std::function}}\label{stdfunction}}

A polymorphic wrapper for any callable (function pointer, lambda,
functor).

\begin{lstlisting}[language={C++}]
#include <functional>

void freeFunc(int) {}

std::function<void(int)> f;
f = freeFunc;
f = [](int x) {};
\end{lstlisting}

\textbf{Cost:} Can incur heap allocation and virtual call overhead.

\begin{center}\rule{0.5\linewidth}{0.5pt}\end{center}

\hypertarget{stdbind}{%
\subsection{\texorpdfstring{3.
\texttt{std::bind}}{3. std::bind}}\label{stdbind}}

Binds arguments to function parameters (Partial Application).

\begin{lstlisting}[language={C++}]
int add(int a, int b) { return a + b; }

// Creates a function taking 1 argument (placeholder _1)
auto add5 = std::bind(add, 5, std::placeholders::_1);
// add5(10) calls add(5, 10)
\end{lstlisting}

\emph{Note: Lambdas largely replace \passthrough{\lstinline!std::bind!}
in modern C++ due to readability and optimization.}

\hypertarget{professional-notes-chapter-52-stdfunction-to-wrap-any-element-that-is-callable}{%
\section{Professional Notes: Chapter 52: std::function: To wrap any
element that is
callable}\label{professional-notes-chapter-52-stdfunction-to-wrap-any-element-that-is-callable}}

Section 52.1: Simple usage Section 52.2: std::function used with
std::bind Section 52.3: Binding std::function to a dierent callable
types Section 52.4: Storing function arguments in std::tuple Section
52.5: std::function with lambda and std::bind Section 52.6:
\passthrough{\lstinline!function!} overhead

\hypertarget{professional-notes-chapter-73-lambdas}{%
\section{Professional Notes: Chapter 73:
Lambdas}\label{professional-notes-chapter-73-lambdas}}

Section 73.1: What is a lambda expression? Section 73.2: Specifying the
return type Section 73.3: Capture by value Section 73.4: Recursive
lambdas Section 73.5: Default capture Section 73.6: Class lambdas and
capture of this Section 73.7: Capture by reference Section 73.8: Generic
lambdas Section 73.9: Using lambdas for inline parameter pack unpacking
Section 73.10: Generalized capture Section 73.11: Conversion to function
pointer Section 73.12: Porting lambda functions to C++03 using functors

\hypertarget{professional-notes-chapter-52-stdfunction-to-wrap-any}{%
\section{Professional Notes: Chapter 52: std::function: To wrap
any}\label{professional-notes-chapter-52-stdfunction-to-wrap-any}}

element that is callable Section 52.1: Simple usage \#include \#include
std::function\textless void(int , const std::string\&)\textgreater{}
myFuncObj; void theFunc(int i, const std::string\& s) \{ std::cout
\textless\textless{} s \textless\textless{} ``:'' \textless\textless{} i
\textless\textless{} std::endl; \} int main(int argc, char *argv{[}{]})
\{ myFuncObj = theFunc; myFuncObj(10, ``hello world''); \} Section 52.2:
std::function used with std::bind Think about a situation where we need
to callback a function with arguments. std::function used with std::bind
gives a very powerful design construct as shown below. class A \{
public: std::function\textless void(int, const
std::string\&)\textgreater{} m\_CbFunc = nullptr; void foo() \{ if
(m\_CbFunc) \{ m\_CbFunc(100, ``event fired''); \} \} \}; class B \{
public: B() \{ auto aFunc = std::bind(\&B::eventHandler, this,
std::placeholders::\_1, std::placeholders::\_2); anObjA.m\_CbFunc =
aFunc; \} void eventHandler(int i, const std::string\& s) \{ std::cout
\textless\textless{} s \textless\textless{} ``:'' \textless\textless{} i
\textless\textless{} std::endl; \} void DoSomethingOnA() \{
anObjA.foo(); \} A anObjA; \}; int main(int argc, char
\emph{argv{[}{]}) \{ B anObjB; anObjB.DoSomethingOnA(); \} Section 52.3:
Binding std::function to a dierent callable types /} * This example show
some ways of using std::function to call * a) C-like function * b)
class-member function * c) operator() * d) lambda function \emph{ }
Function call can be made: * a) with right arguments * b) argumens with
different order, types and count */ \#include \#include \#include
\#include using std::cout; using std::endl; using namespace
std::placeholders; // simple function to be called double foo\_fn(int x,
float y, double z) \{ double res = x + y + z; std::cout
\textless\textless{} ``foo\_fn called with arguments:''
\textless\textless{} x \textless\textless{} ``,'' \textless\textless{} y
\textless\textless{} ``,'' \textless\textless{} z \textless\textless{} "
result is : " \textless\textless{} res \textless\textless{} std::endl;
return res; \} // structure with member function to call struct
foo\_struct \{ // member function to call double foo\_fn(int x, float y,
double z) \{ double res = x + y + z; std::cout \textless\textless{}
``foo\_struct::foo\_fn called with arguments:'' \textless\textless{} x
\textless\textless{} ``,'' \textless\textless{} y \textless\textless{}
``,'' \textless\textless{} z \textless\textless{} " result is : "
\textless\textless{} res \textless\textless{} std::endl; return res; \}
// this member function has different signature - but it can be used too
// please not that argument order is changed too double foo\_fn\_4(int
x, double z, float y, long xx) \{ double res = x + y + z + xx; std::cout
\textless\textless{} ``foo\_struct::foo\_fn\_4 called with arguments:''
 \textless\textless{} x \textless\textless{} ``,'' \textless\textless{}
z \textless\textless{} ``,'' \textless\textless{} y \textless\textless{}
``,'' \textless\textless{} xx \textless\textless{} " result is : "
\textless\textless{} res \textless\textless{} std::endl; return res; \}
// overloaded operator() makes whole object to be callable double
operator()(int x, float y, double z) \{ double res = x + y + z;
std::cout \textless\textless{} ``foo\_struct::operator() called with
arguments:'' \textless\textless{} x \textless\textless{} ``,''
\textless\textless{} y \textless\textless{} ``,'' \textless\textless{} z
\textless\textless{} " result is : " \textless\textless{} res
\textless\textless{} std::endl; return res; \} \}; int main(void) \{ //
typedefs using function\_type = std::function\textless double(int,
float, double)\textgreater; // foo\_struct instance foo\_struct fs; //
here we will store all binded functions std::vector bindings; // var \#1
- you can use simple function function\_type var1 = foo\_fn;
bindings.push\_back(var1); // var \#2 - you can use member function
function\_type var2 = std::bind(\&foo\_struct::foo\_fn, fs, \_1, \_2,
\_3); bindings.push\_back(var2); // var \#3 - you can use member
function with different signature // foo\_fn\_4 has different count of
arguments and types function\_type var3 =
std::bind(\&foo\_struct::foo\_fn\_4, fs, \_1, \_3, \_2, 0l);
bindings.push\_back(var3); // var \#4 - you can use object with
overloaded operator() function\_type var4 = fs;
bindings.push\_back(var4); // var \#5 - you can use lambda function
function\_type var5 = \href{int\%20x,\%20float\%20y,\%20double\%20z}{}
\{ double res = x + y + z; std::cout \textless\textless{} ``lambda
called with arguments:'' \textless\textless{} x \textless\textless{}
``,'' \textless\textless{} y \textless\textless{} ``,''
\textless\textless{} z \textless\textless{} " result is : "
\textless\textless{} res \textless\textless{} std::endl; return res; \};
bindings.push\_back(var5); std::cout \textless\textless{} ``Test stored
functions with arguments: x = 1, y = 2, z = 3'' \textless\textless{}
std::endl; for (auto f : bindings)  f(1, 2, 3); \} Live Output: Test
stored functions with arguments: x = 1, y = 2, z = 3 foo\_fn called with
arguments: 1, 2, 3 result is : 6 foo\_struct::foo\_fn called with
arguments: 1, 2, 3 result is : 6 foo\_struct::foo\_fn\_4 called with
arguments: 1, 3, 2, 0 result is : 6 foo\_struct::operator() called with
arguments: 1, 2, 3 result is : 6 lambda called with arguments: 1, 2, 3
result is : 6 Section 52.4: Storing function arguments in std::tuple
Some programs need so store arguments for future calling of some
function. This example shows how to call any function with arguments
stored in std::tuple \#include \#include \#include \#include // simple
function to be called double foo\_fn(int x, float y, double z) \{ double
res = x + y + z; std::cout \textless\textless{} ``foo\_fn called. x =''
\textless\textless{} x \textless\textless{} " y = " \textless\textless{}
y \textless\textless{} " z = " \textless\textless{} z
\textless\textless{} " res=" \textless\textless{} res; return res; \} //
helpers for tuple unrolling template\textless int \ldots\textgreater{}
struct seq \{\}; template\textless int N, int \ldots S\textgreater{}
struct gens : gens\textless N-1, N-1, S\ldots\textgreater{} \{\};
template\textless int \ldots S\textgreater{} struct gens\textless0,
S\ldots\textgreater\{ typedef seq\textless S\ldots\textgreater{} type;
\}; // invocation helper template\textless typename FN, typename P, int
\ldots S\textgreater{} double call\_fn\_internal(const FN\& fn, const
P\& params, const seq\textless S\ldots\textgreater) \{ return
fn(std::get(params) \ldots); \} // call function with arguments stored
in std::tuple template\textless typename Ret, typename
\ldots Args\textgreater{} Ret call\_fn(const
std::function\textless Ret(Args\ldots)\textgreater\& fn, const
std::tuple\textless Args\ldots\textgreater\& params) \{ return
call\_fn\_internal(fn, params, typename
gens\textless sizeof\ldots(Args)\textgreater::type()); \} int main(void)
\{ // arguments std::tuple\textless int, float, double\textgreater{} t =
std::make\_tuple(1, 5, 10); // function to call 
std::function\textless double(int, float, double)\textgreater{} fn =
foo\_fn; // invoke a function with stored arguments call\_fn(fn, t); \}
Live Output: foo\_fn called. x = 1 y = 5 z = 10 res=16 Section 52.5:
std::function with lambda and std::bind \#include \#include using
std::placeholders::\_1; // to be used in std::bind example int
stdf\_foobar (int x, std::function\textless int(int)\textgreater{} moo)
\{ return x + moo(x); // std::function moo called \} int foo (int x) \{
return 2+x; \} int foo\_2 (int x, int y) \{ return 9\emph{x + y; \} int
main() \{ int a = 2; /} Function pointers \emph{/ std::cout
\textless\textless{} stdf\_foobar(a, \&foo) \textless\textless{}
std::endl; // 6 ( 2 + (2+2) ) // can also be: stdf\_foobar(2, foo) /}
Lambda expressions \emph{/ /} An unnamed closure from a lambda
expression can be * stored in a std::function object: \emph{/ int
capture\_value = 3; std::cout \textless\textless{} stdf\_foobar(a,
\href{int\%20param}{capture\_value} -\textgreater{} int \{ return 7 +
capture\_value } param; \}) \textless\textless{} std::endl; // result:
15 == value + (7 * capture\_value * value) == 2 + (7 + 3 * 2) /*
std::bind expressions \emph{/ /} The result of a std::bind expression
can be passed. * For example by binding parameters to a function pointer
call: */ int b = stdf\_foobar(a, std::bind(foo\_2, \_1, 3)); std::cout
\textless\textless{} b \textless\textless{} std::endl; // b == 23 == 2 +
( 9*2 + 3 ) int c = stdf\_foobar(a, std::bind(foo\_2, 5, \_1));
std::cout \textless\textless{} c \textless\textless{} std::endl; // c ==
49 == 2 + ( 9*5 + 2 ) return 0; \} Section 52.6:
\passthrough{\lstinline!function!} overhead std::function can cause
signicant overhead. Because std::function has {[}value
semantics{]}{[}1{]}, it must copy or move the given callable into
itself. But since it can take callables of an arbitrary type, it will
frequently have to allocate memory dynamically to do this. Some function
implementations have so-called ``small object optimization'', where
small types (like function pointers, member pointers, or functors with
very little state) will be stored directly in the function object. But
even this only works if the type is noexcept move constructible.
Furthermore, the C++ standard does not require that all implementations
provide one. Consider the following: //Header file using MyPredicate =
std::function\textless bool(const MyValue \&, const MyValue
\&)\textgreater; void SortMyContainer(MyContainer \&C, const MyPredicate
\&pred); //Source file void SortMyContainer(MyContainer \&C, const
MyPredicate \&pred) \{ std::sort(C.begin(), C.end(), pred); \} A
template parameter would be the preferred solution for SortMyContainer,
but let us assume that this is not possible or desirable for whatever
reason. SortMyContainer does not need to store pred beyond its own call.
And yet, pred may well allocate memory if the functor given to it is of
some non-trivial size. function allocates memory because it needs
something to copy/move into; function takes ownership of the callable it
is given. But SortMyContainer does not need to own the callable; it's
just referencing it. So using function here is overkill; it may be
ecient, but it may not. There is no standard library function type that
merely references a callable. So an alternate solution will have to be
found, or you can choose to live with the overhead. Also, function has
no eective means to control where the memory allocations for the object
come from. Yes, it has constructors that take an allocator, but {[}many
implementations do not implement them correctly\ldots{} or even at
all{]}{[}2{]}. Version C++17 The function constructors that take an
allocator no longer are part of the type. Therefore, there is no way to
manage the allocation. Calling a function is also slower than calling
the contents directly. Since any function instance could hold any
callable, the call through a function must be indirect. The overhead of
calling function is on the order of a virtual function call.

\hypertarget{professional-notes-chapter-73-lambdas-1}{%
\section{Professional Notes: Chapter 73:
Lambdas}\label{professional-notes-chapter-73-lambdas-1}}

Parameter default-capture Details Species how all non-listed variables
are captured. Can be = (capture by value) or \& (capture by reference).
If omitted, non-listed variables are inaccessible within the
lambda-body. The default- capture must precede the capture-list.
capture-list Species how local variables are made accessible within the
lambda-body. Variables without prex are captured by value. Variables
prexed with \& are captured by reference. Within a class method, this
can be used to make all its members accessible by reference. Non-listed
variables are inaccessible, unless the list is preceded by a
default-capture. argument-list Species the arguments of the lambda
function. mutable (optional) Normally variables captured by value are
const. Specifying mutable makes them non- const. Changes to those
variables are retained between calls. throw-specication (optional)
Species the exception throwing behavior of the lambda function. For
example: noexcept or throw(std::exception). attributes (optional) Any
attributes for the lambda function. For example, if the lambda-body
always throws an exception then {[}{[}noreturn{]}{]} can be used.
-\textgreater{} return-type (optional) Species the return type of the
lambda function. Required when the return type cannot be determined by
the compiler. lambda-body A code block containing the implementation of
the lambda function. Section 73.1: What is a lambda expression? A lambda
expression provides a concise way to create simple function objects. A
lambda expression is a prvalue whose result object is called closure
object, which behaves like a function object. The name `lambda
expression' originates from lambda calculus, which is a mathematical
formalism invented in the 1930s by Alonzo Church to investigate
questions about logic and computability. Lambda calculus formed the
basis of LISP, a functional programming language. Compared to lambda
calculus and LISP, C++ lambda expressions share the properties of being
unnamed, and to capture variables from the surrounding context, but they
lack the ability to operate on and return functions. A lambda expression
is often used as an argument to functions that take a callable object.
That can be simpler than creating a named function, which would be only
used when passed as the argument. In such cases, lambda expressions are
generally preferred because they allow dening the function objects
inline. A lambda consists typically of three parts: a capture list
{[}{]}, an optional parameter list () and a body \{\}, all of which can
be empty: \href{}{} // An empty lambda, which does and returns nothing
Capture list {[}{]} is the capture list. By default, variables of the
enclosing scope cannot be accessed by a lambda. Capturing a variable
makes it accessible inside the lambda, either as a copy or as a
reference. Captured variables become a part of the lambda; in contrast
to function arguments, they do not have to be passed when calling the
lambda. int a = 0; // Define an integer variable auto f = \href{}{} \{
return a\emph{9; \}; // Error: `a' cannot be accessed auto f =
\href{}{a} \{ return a}9; \}; // OK, `a' is ``captured'' by value auto f
= \href{}{\&a} \{ return a++; \}; // OK, `a' is ``captured'' by
reference // Note: It is the responsibility of the programmer // to
ensure that a is not destroyed before the
